\documentclass[aspectratio=169,10pt,t]{beamer}
\usepackage{ac1}
\aulainfo{34 e 35}{Apresentação dos Projetos}{Demonstração, explicação e critérios de avaliação}

\begin{document}
\capa

\begin{frame}{Como serão as apresentações}
  \begin{itemize}
    \item Cada grupo apresenta seu projeto para a turma: \textbf{10 a 12 minutos} + 3 minutos de perguntas
    \item Todos os integrantes falam
    \item Ordem definida por sorteio no início da aula 34
    \item \textbf{Entrega do código} (pasta compactada com instruções) antes da apresentação
  \end{itemize}
  \begin{alertblock}{Atenção}
    Teste a demonstração \emph{no computador que será usado}. Tenha um plano B (capturas de tela ou vídeo curto) caso algo falhe.
  \end{alertblock}
\end{frame}

\begin{frame}{Roteiro sugerido (10 minutos)}
  \centering\small
  \begin{tabular}{p{1.3cm}p{7.4cm}p{2.6cm}}
    \toprule
    \textbf{Tempo} & \textbf{Conteúdo} & \textbf{Quem}\\
    \midrule
    1 min & O problema: qual é, para quem, por que importa & integrante A\\
    2 min & Visão geral: funcionalidades e estrutura do programa (módulos/funções) & integrante B\\
    4 min & \textbf{Demonstração ao vivo} com dados reais & integrante A ou C\\
    2 min & Decisões de projeto e desafios enfrentados (o que deu errado e como resolveram) & integrante C\\
    1 min & Conclusões: limitações e próximos passos & todos\\
    \bottomrule
  \end{tabular}
\end{frame}

\begin{frame}{Critérios de avaliação}
  \begin{columns}[T]
    \column{0.5\textwidth}
    \begin{block}{Funcionalidade}
      O programa atende aos requisitos especificados? Funciona com dados diferentes dos do exemplo? Trata entradas inválidas?
    \end{block}
    \begin{block}{Organização do código}
      O código é bem estruturado e fácil de entender? Nomes claros, docstrings, sem repetição?
    \end{block}
    \column{0.5\textwidth}
    \begin{block}{Modularização}
      O código está organizado em funções e módulos de forma lógica, com cada parte tendo uma responsabilidade?
    \end{block}
    \begin{block}{Clareza na explicação}
      Os alunos explicam claramente as decisões de design e a implementação? Respondem às perguntas?
    \end{block}
  \end{columns}
\end{frame}

\begin{frame}{Pontos extras e pontos de atenção}
  \begin{columns}[T]
    \column{0.5\textwidth}
    \begin{exampleblock}{Valorizamos}
      \begin{itemize}\small
        \item Testes automáticos
        \item Gráficos bem feitos e interpretados
        \item Tratamento cuidadoso de erros
        \item Uso adequado de NumPy/Pandas/Matplotlib
        \item Instruções de uso claras (\texttt{LEIAME})
      \end{itemize}
    \end{exampleblock}
    \column{0.5\textwidth}
    \begin{alertblock}{Perdem pontos}
      \begin{itemize}\small
        \item Código que só um integrante entende
        \item Programa que ``quebra'' em entradas simples
        \item Funções gigantes e repetição
        \item Apresentação que só lê os slides
        \item Estourar o tempo
      \end{itemize}
    \end{alertblock}
  \end{columns}
\end{frame}

\begin{frame}{Perguntas que o professor poderá fazer}
  \begin{itemize}
    \item Por que escolheram \emph{esta} estrutura de dados (lista, dicionário, DataFrame)?
    \item O que acontece se o arquivo não existir? E se o usuário digitar texto no lugar de número?
    \item Mostre a função mais importante e explique linha a linha
    \item Como vocês testaram? Qual foi o bug mais difícil?
    \item Onde há variáveis globais? Como removê-las?
    \item O que faria diferente se começasse hoje? O que falta para uso real?
  \end{itemize}
\end{frame}

\begin{frame}{Dicas para uma boa apresentação}
  \begin{itemize}
    \item \textbf{Comece pelo problema}, não pelo código
    \item Poucos slides, poucas palavras; mostre o \emph{programa rodando}
    \item Use entradas \textbf{realistas}; mostre também um erro sendo tratado
    \item Mostre trechos curtos de código (letra grande!) para ilustrar uma decisão
    \item Treine antes e cronometre
    \item Assuma as limitações: ``ainda não faz X'' é melhor que ser surpreendido
  \end{itemize}
  \begin{exampleblock}{Material de apoio}
    Cap.~9 do livro e o site \href{\siteurl}{\texttt{\siteurl}}
  \end{exampleblock}
\end{frame}

\begin{frame}{Encerramento da disciplina}
  \begin{itemize}
    \item O que aprendemos: pensar algoritmicamente, programar em Python, organizar dados, analisar e visualizar
    \item Isto é o \textbf{começo}: as próximas disciplinas e a vida profissional exigirão programar cada vez mais
    \item Continue praticando: exercícios adicionais do site, projetos pessoais, dados da sua área
  \end{itemize}
  \vspace{1ex}
  \begin{block}{Para continuar}
    Estruturas de dados e algoritmos \(\cdot\) Cálculo numérico \(\cdot\) Estatística com Python \(\cdot\) Banco de dados \(\cdot\) Aprendizado de máquina
  \end{block}
  \vspace{1ex}
  \centering\large Obrigado!
\end{frame}

\end{document}
